%% A．導入と背景（スライド 2〜6．題はスライド 1）
%% 専門家に向けた版（2026-11-10 研究者「思いっきり内容変えていい．もっと専門家向けに」）．英語は後で訳す（\EN の印）
\setband{\T{導入と背景}{\EN}}

% ---------------------------------------------------------------------
\begin{frame}{\T{課題：新しいコーパスで，感情ごとに K 本のラベルから MER の分類器を作る}{\EN}}
\begin{columns}[T]
\begin{column}{0.47\textwidth}
\small
\T{%
\begin{itemize}
\item \textbf{入力}：発話ごとの特徴 $x = (x^{\mathrm{A}}, x^{\mathrm{V}}, x^{\mathrm{T}})$（声・顔・言葉．凍結したエンコーダの出力）
\item \textbf{出力}：$\hat{y} \in$ \{平静，喜び，悲しみ，怒り\}
\item \textbf{ソース}：ほかのコーパス．全発話にラベルがある
\item \textbf{ターゲット}：感情ごとに K 本のラベル（K ∈ \{0, 1, 5, 10, 25\}）と，ラベル無しの発話
\item \textbf{評価}：ターゲットのテストの話者での UAR．テストの話者は学習に一切使わない
\end{itemize}}{\EN}
\end{column}
\begin{column}{0.51\textwidth}
\centering
\begin{adjustbox}{max width=\linewidth}
\begin{tikzpicture}[x=1cm, y=1cm]
  \foreach \n/\y in {\T{ソース 1}{\EN}/1.3, \T{ソース 2}{\EN}/0.2, \T{ソース 3}{\EN}/-0.9} {
    \node[srcbox, minimum width=1.8cm, minimum height=0.8cm] at (0,\y) {\n\\[-2pt]{\tiny\T{全発話にラベル}{\EN}}};
  }
  \node[draw=mDarkTeal!70, fill=cTgt!50, rounded corners=3pt, minimum width=3.6cm, minimum height=3.5cm] (tgt) at (3.9,0.2) {};
  \node[font=\scriptsize\bfseries, color=mDarkTeal] at (3.9,2.15) {\T{ターゲット（新しいコーパス）}{\EN}};
  \foreach \c [count=\r from 0] in {eNeu,eHap,eSad,eAng} {\foreach \i in {0,...,4} {\node[utt, fill=\c, draw=\c, minimum size=2.4mm] at (2.55+0.3*\i,1.55-0.3*\r) {};}}
  \node[font=\tiny, anchor=west, align=left] at (4.1,1.1) {\T{K 本 × 4 感情}{\EN}\\\T{（図は K = 5）}{\EN}};
  \foreach \r in {0,1} {\foreach \i in {0,...,9} {\node[utt, minimum size=2.4mm] at (2.55+0.3*\i,0.0-0.3*\r) {};}}
  \node[font=\tiny, anchor=west] at (2.4,-0.6) {\T{ラベル無しの発話（候補）}{\EN}};
  \node[testbox, minimum width=3.1cm, minimum height=0.5cm, font=\tiny] at (3.9,-1.15) {\T{テスト（別の話者．ラベルは評価だけに使う）}{\EN}};
  \draw[arr] (1.0,0.2) -- (2.05,0.2);
\end{tikzpicture}
\end{adjustbox}
\end{column}
\end{columns}
\fnote{\T{UAR（unweighted average recall）：感情ごとの再現率の平均．感情の割合の偏りに左右されず，でたらめに答えると 25\%．この設定は，少数ラベルの半教師ありドメイン適応にあたる}{\EN}}
\end{frame}

% ---------------------------------------------------------------------
\begin{frame}{\T{コーパスが変わると成績が大きく落ちる（クロスコーパスの劣化）}{\EN}}
\centering
\begin{tikzpicture}
\begin{axis}[width=0.78\linewidth, height=4.5cm, ybar, bar width=14pt, ymin=0, ymax=100, enlarge x limits=0.18,
  xtick={1,2,3,4}, xticklabel style={align=center, font=\small},
  xticklabels={{AFEW\\\T{映画}{\EN}}, {IEMOCAP\\\T{即興の対話}{\EN}}, {MELD\\\T{ドラマ}{\EN}}, {MOSEI\\\T{YouTube}{\EN}}},
  ylabel={UAR (\%)}, ytick={0,25,50,75,100},
  legend style={at={(0.5,1.03)}, anchor=south, font=\small, legend columns=2, /tikz/every even column/.append style={column sep=1.2em}},
  nodes near coords={\pgfmathprintnumber[fixed, precision=0]{\pgfplotspointmeta}}, every node near coord/.append style={font=\scriptsize},
  extra y ticks={25}, extra y tick style={grid=major, major grid style={mMuted, dashed}}, extra y tick labels={}]
\addplot[fill=cSame, draw=cSame] table[x=ci, y=upper, col sep=comma, comment chars={\#}]{figures/intro_drop.csv};
\addplot[fill=cOther, draw=cOther] table[x=ci, y=start, col sep=comma, comment chars={\#}]{figures/intro_drop.csv};
\legend{\T{同じコーパスで学習（全ラベル）}{\EN}, \T{ほかの 3 コーパスで学習（ターゲットのラベル 0 本）}{\EN}}
\end{axis}
\end{tikzpicture}
\msg{\T{融合の分類器（声・顔・言葉の特徴をつないだもの）でも \N{igapmin}〜\N{igapmax} ポイント下がる}{\EN}}
\fnote{\T{ほかの 3 コーパスで学習：コーパスごとに学んだ線形分類器の重みと切片の平均．同じコーパスで学習：テストに使わない話者の全ラベルで学び，切片だけをテストの発話のラベルで解き直した値（目安の上限）．コーパスごとに 5 つに分けた話者の 1 つをテストにする．LOCO 設定（leave-one-corpus-out：4 コーパスのうち 1 つをターゲット，残りをソースにする）の評価の回の値．出どころ：figures/intro\_drop.csv}{\EN}}
\end{frame}

% ---------------------------------------------------------------------
\begin{frame}{\T{使用したコーパス：6 コーパス 7 群，4 感情}{\EN}}
\centering
\scriptsize
\renewcommand{\arraystretch}{1.18}
\setlength{\tabcolsep}{4pt}
\begin{tabular}{@{}llrrrrrrcccc@{}}
\toprule
\T{群}{\EN} & \T{収録}{\EN} & \T{話者群}{\EN} & \textcolor{eNeu}{\T{平静}{\EN}} & \textcolor{eHap}{\T{喜び}{\EN}} & \textcolor{eSad}{\T{悲しみ}{\EN}} & \textcolor{eAng}{\T{怒り}{\EN}} & \T{計}{\EN} & \textcolor{kA}{\T{声}{\EN}} & \textcolor{kV}{\T{顔}{\EN}} & \textcolor{kT}{\T{言葉}{\EN}} & LOCO \\
\midrule
IEMOCAP \T{即興}{\EN} & \T{演技者 10 人の即興の対話}{\EN} & \N{cspkiemoimpro} & \N{cneuiemoimpro} & \N{chapiemoimpro} & \N{csadiemoimpro} & \N{cangiemoimpro} & \N{ctotiemoimpro} & \checkmark & \checkmark & \checkmark & \checkmark \\
IEMOCAP \T{台本}{\EN} & \T{同じ 10 人が台本を演じる対話}{\EN} & \N{cspkiemoscript} & \N{cneuiemoscript} & \N{chapiemoscript} & \N{csadiemoscript} & \N{cangiemoscript} & \N{ctotiemoscript} & \checkmark & \checkmark & \checkmark & \\
MOSEI & \T{YouTube の独白}{\EN} & \N{cspkmosei} & \N{cneumosei} & \N{chapmosei} & \N{csadmosei} & \N{cangmosei} & \N{ctotmosei} & \checkmark & \checkmark & \checkmark & \checkmark \\
MELD & \T{ドラマ（Friends）の多人数の会話}{\EN} & \N{cspkmeld} & \N{cneumeld} & \N{chapmeld} & \N{csadmeld} & \N{cangmeld} & \N{ctotmeld} & \checkmark & \checkmark & \checkmark & \checkmark \\
AFEW & \T{映画の切り抜き}{\EN} & \N{cspkafew} & \N{cneuafew} & \N{chapafew} & \N{csadafew} & \N{cangafew} & \N{ctotafew} & \checkmark & \checkmark & \checkmark$^{*}$ & \checkmark \\
RAVDESS & \T{演技者が決まった文を読む}{\EN} & \N{cspkravdess} & \N{cneuravdess} & \N{chapravdess} & \N{csadravdess} & \N{cangravdess} & \N{ctotravdess} & \checkmark & & & \\
CREMA-D & \T{演技者が決まった文を読む}{\EN} & \N{cspkcremad} & \N{cneucremad} & \N{chapcremad} & \N{csadcremad} & \N{cangcremad} & \N{ctotcremad} & \checkmark & \checkmark & & \\
\bottomrule
\end{tabular}
\msg{\T{収録のしかた，感情の割合，ラベルの決め方がコーパスごとに違う}{\EN}}
\fnote{\T{件数は音声（WavLM）の特徴が取れた発話（data/README.md．顔と言葉では数本〜数十本違う）．$^{*}$ AFEW の言葉は Whisper large-v3 の書き起こし．MOSEI の話者群の多くは動画 ID．ラベル：IEMOCAP の喜びは happy と excited を合わせ，frustrated は怒りに入れない．MOSEI の平静は「どの感情の強さも 0.5 未満」で決まる．LOCO：leave-one-corpus-out．4 コーパスのうち 1 つをターゲット，残りをソースにする．LOCO では 3 モダリティがそろう発話だけを使う}{\EN}}
\end{frame}

% ---------------------------------------------------------------------
\begin{frame}{\T{特徴：凍結した 3 つのエンコーダ．分類器：線形（多クラスのロジスティック回帰）}{\EN}}
\centering
\footnotesize
\renewcommand{\arraystretch}{1.15}
\begin{tabular}{@{}llllr@{}}
\toprule
& \T{エンコーダ}{\EN} & \T{事前学習}{\EN} & \T{発話の特徴の作り方}{\EN} & \T{次元}{\EN} \\
\midrule
\textcolor{kA}{\textbf{\T{声}{\EN}}} & WavLM-Large & \T{自己教師あり（音声）}{\EN} & \T{最終層を時間で平均}{\EN} & 1024 \\
\textcolor{kV}{\textbf{\T{顔}{\EN}}} & ResNet-50 & \T{VGGFace2 → AffectNet で微調整}{\EN} & \T{顔のフレームの特徴を平均}{\EN} & 512 \\
\textcolor{kT}{\textbf{\T{言葉}{\EN}}} & RoBERTa-base & \T{自己教師あり（テキスト）}{\EN} & \T{最終層を単語で平均}{\EN} & 768 \\
\textbf{\T{融合}{\EN}} & — & — & \T{3 つの特徴をつなぐ}{\EN} & 2,304 \\
\bottomrule
\end{tabular}

\vspace{0.3em}
\begin{tikzpicture}[x=1cm, y=1cm, every node/.append style={font=\footnotesize}]
  \node[big, fill=cFixed!40, draw=mMuted, minimum width=2.6cm] (e) at (0,0) {\T{エンコーダ}{\EN}\\[-2pt]{\scriptsize\T{凍結}{\EN}}};
  \node[big, draw=mDarkTeal!60, minimum width=2.6cm] (z) at (3.6,0) {\T{コーパスごとに標準化}{\EN}\\[-2pt]{\scriptsize $x$}};
  \node[big, fill=cLearn!10, draw=cLearn, minimum width=3.2cm] (c) at (7.6,0) {\T{線形分類器}{\EN}\\[-2pt]{\scriptsize $\hat{y} = \arg\max_k\, (W x + b)_k$}};
  \draw[arr] (e) -- (z); \draw[arr] (z) -- (c);
\end{tikzpicture}
\msg{\T{凍結は解析のための設定：誤りを切片 $b$ と重み $W$ に分けて測る}{\EN}}
\fnote{\T{顔は Ryumina らの静的な表情モデル．確かめに第 2 のエンコーダ（HuBERT-Large，CLIP ViT-L/14，DeBERTa-v3-base）も使う．凍結＋軽い分類器は SUPERB・EmoBox・MERBench の標準の評価の形}{\EN}}
\end{frame}
